% GENERATED FILE. Edit canonical lessons, then run npm run generate:books.
% canonical-source-hash: fnv1a64:56add179574fd9ea
% canonical-lessons: TE-C171-suryastamayam, TE-C171-pratiroju, TE-C171-sadharananga, TE-C171-aruduga, TE-C171-nirantaram

\chapter{Sunset, Every day, Usually, Rarely, Constantly}
\label{ch:te-sunset-every-day-usually-rarely-constantly}
\hlchaptermodality{171}

\begin{chapteropening}
\textbf{By the end of this chapter:} \emph{I can say sunset, every day, usually, rarely and constantly.}

Complete the last lesson of chapter 171: I can say sunset, every day, usually, rarely and constantly.
\end{chapteropening}

\section[\texorpdfstring{\te{సూర్యాస్తమయం} (sūryāstamayaṁ)}{సూర్యాస్తమయం (sūryāstamayaṁ)}]{\te{సూర్యాస్తమయం} (sūryāstamayaṁ) — sunset}

\label{lesson:TE-C171-suryastamayam}



\begin{quote}
Before the new one: say the Telugu for a date, then the Telugu for a weekend.
\end{quote}

\subsection*{You'll want to know: \te{సూర్యాస్తమయం}}
\textbf{\te{సూర్యాస్తమయం}} — \emph{sūryāstamayaṁ} — \textquotedblleft{}sunset\textquotedblright{}.

\begin{grammarlens}[title={What you've built}]
One more word for this chapter.
\end{grammarlens}

\subsection*{Your turn}
\begin{itemize}
\raggedright
  \item \emph{Say it:} \emph{sūryāstamayaṁ}
  \item \emph{Say it:} \emph{sūryāstamayaṁ}, once more
  \item \emph{Say it:} say \emph{sūryāstamayaṁ}
  \item \emph{Recall:} say the Telugu for a date, then the Telugu for a weekend, then say \emph{sūryāstamayaṁ} again
\end{itemize}

\begin{tcolorbox}[breakable,colback=teal!4,colframe=teal!35!black,title={Before you move on}]
What does \te{సూర్యాస్తమయం} mean? (\textquotedblleft{}Sunset\textquotedblright{}.) Say it once more.
\end{tcolorbox}
\section[\texorpdfstring{\te{ప్రతిరోజు} (pratirōju)}{ప్రతిరోజు (pratirōju)}]{\te{ప్రతిరోజు} (pratirōju) — every day}

\label{lesson:TE-C171-pratiroju}



\begin{quote}
Before the new one: say the Telugu for a weekend, then the Telugu for sunset.
\end{quote}

\subsection*{You'll want to know: \te{ప్రతిరోజు}}
\textbf{\te{ప్రతిరోజు}} — \emph{pratirōju} — \textquotedblleft{}every day\textquotedblright{}.

\begin{grammarlens}[title={What you've built}]
One more word for this chapter.
\end{grammarlens}

\subsection*{Your turn}
\begin{itemize}
\raggedright
  \item \emph{Say it:} \emph{pratirōju}
  \item \emph{Say it:} \emph{pratirōju}, once more
  \item \emph{Say it:} say \emph{pratirōju}
  \item \emph{Recall:} say the Telugu for a weekend, then the Telugu for sunset, then say \emph{pratirōju} again
\end{itemize}

\begin{tcolorbox}[breakable,colback=teal!4,colframe=teal!35!black,title={Before you move on}]
What does \te{ప్రతిరోజు} mean? (\textquotedblleft{}Every day\textquotedblright{}.) Say it once more.
\end{tcolorbox}
\section[\texorpdfstring{\te{సాధారణంగా} (sādhāraṇaṅgā)}{సాధారణంగా (sādhāraṇaṅgā)}]{\te{సాధారణంగా} (sādhāraṇaṅgā) — usually, generally}

\label{lesson:TE-C171-sadharananga}



\begin{quote}
Before the new one: say the Telugu for sunset, then the Telugu for every day.
\end{quote}

\subsection*{You'll want to know: \te{సాధారణంగా}}
\textbf{\te{సాధారణంగా}} — \emph{sādhāraṇaṅgā} — \textquotedblleft{}usually, generally\textquotedblright{}.

\begin{grammarlens}[title={What you've built}]
One more word for this chapter.
\end{grammarlens}

\subsection*{Your turn}
\begin{itemize}
\raggedright
  \item \emph{Say it:} \emph{sādhāraṇaṅgā}
  \item \emph{Say it:} \emph{sādhāraṇaṅgā}, once more
  \item \emph{Say it:} say \emph{sādhāraṇaṅgā}
  \item \emph{Recall:} say the Telugu for sunset, then the Telugu for every day, then say \emph{sādhāraṇaṅgā} again
\end{itemize}

\begin{tcolorbox}[breakable,colback=teal!4,colframe=teal!35!black,title={Before you move on}]
What does \te{సాధారణంగా} mean? (\textquotedblleft{}Usually, generally\textquotedblright{}.) Say it once more.
\end{tcolorbox}
\section[\texorpdfstring{\te{అరుదుగా} (arudugā)}{అరుదుగా (arudugā)}]{\te{అరుదుగా} (arudugā) — rarely}

\label{lesson:TE-C171-aruduga}



\begin{quote}
Before the new one: say the Telugu for every day, then the Telugu for usually.
\end{quote}

\subsection*{You'll want to know: \te{అరుదుగా}}
\textbf{\te{అరుదుగా}} — \emph{arudugā} — \textquotedblleft{}rarely\textquotedblright{}.

\begin{grammarlens}[title={What you've built}]
One more word for this chapter.
\end{grammarlens}

\subsection*{Your turn}
\begin{itemize}
\raggedright
  \item \emph{Say it:} \emph{arudugā}
  \item \emph{Say it:} \emph{arudugā}, once more
  \item \emph{Say it:} say \emph{arudugā}
  \item \emph{Recall:} say the Telugu for every day, then the Telugu for usually, then say \emph{arudugā} again
\end{itemize}

\begin{tcolorbox}[breakable,colback=teal!4,colframe=teal!35!black,title={Before you move on}]
What does \te{అరుదుగా} mean? (\textquotedblleft{}Rarely\textquotedblright{}.) Say it once more.
\end{tcolorbox}
\section[\texorpdfstring{\te{నిరంతరం} (nirantaraṁ)}{నిరంతరం (nirantaraṁ)}]{\te{నిరంతరం} (nirantaraṁ) — constantly, continuously}

\label{lesson:TE-C171-nirantaram}



\begin{quote}
Before the new one: say the Telugu for usually, then the Telugu for rarely.
\end{quote}

\subsection*{You'll want to know: \te{నిరంతరం}}
\textbf{\te{నిరంతరం}} — \emph{nirantaraṁ} — \textquotedblleft{}constantly, continuously\textquotedblright{}. \emph{Read it:} two words again — \textbf{\te{ఛాయ}} (\emph{chāya}, shade) is written with \te{ఛ}, and \textbf{\te{ఝరి}} (\emph{jhari}, a stream) with \te{ఝ}

\begin{grammarlens}[title={What you've built}]
One more word for this chapter.
\end{grammarlens}

\subsection*{Your turn}
\begin{itemize}
\raggedright
  \item \emph{Say it:} \emph{nirantaraṁ}
  \item \emph{Say it:} \emph{nirantaraṁ}, once more
  \item \emph{Say it:} say \emph{nirantaraṁ}
  \item \emph{Recall:} say the Telugu for usually, then the Telugu for rarely, then say \emph{nirantaraṁ} again
\end{itemize}

\begin{tcolorbox}[breakable,colback=teal!4,colframe=teal!35!black,title={Before you move on}]
What does \te{నిరంతరం} mean? (\textquotedblleft{}Constantly, continuously\textquotedblright{}.) Say it once more.
\end{tcolorbox}
